\documentclass[conference]{IEEEtran}
\IEEEoverridecommandlockouts
% The preceding line is only needed to identify funding in the first footnote. If that is unneeded, please comment it out.
\usepackage{cite}
\usepackage{amsmath,amssymb,amsfonts}
\usepackage{algorithmic}
\usepackage{graphicx}
\usepackage{textcomp}
\usepackage{xcolor}
\def\BibTeX{{\rm B\kern-.05em{\sc i\kern-.025em b}\kern-.08em
    T\kern-.1667em\lower.7ex\hbox{E}\kern-.125emX}}

\begin{document}

\title{MedPredictX: An AI-Powered Healthcare Triage and Management System\\
}

\author{\IEEEauthorblockN{1\textsuperscript{st} Given Name Surname}
\IEEEauthorblockA{\textit{dept. name of organization (of Aff.)} \\
\textit{name of organization (of Aff.)}\\
City, Country \\
email address or ORCID}
}

\maketitle

\begin{abstract}
The rapid advancement of artificial intelligence and machine learning has opened new paradigms in the healthcare sector, particularly in the domain of automated medical triage and preliminary diagnosis. This paper presents MedPredictX, a comprehensive, full-stack healthcare application designed to seamlessly bridge the gap between patients and medical professionals. MedPredictX integrates a robust Django-based backend architecture with a modern React frontend, supplemented by state-of-the-art AI capabilities. By leveraging Large Language Models (LLMs) via the Groq API for intelligent medical recommendations and a specialized local Machine Learning model (Random Forest) for precise disease prediction based on 132 distinct symptoms, the system provides instantaneous, reliable diagnostic insights. Furthermore, the platform features a complete suite for patient management, including appointment scheduling, medical record tracking, and secure JWT-based authentication. This report details the system architecture, methodology, implementation nuances, and the results of integrating AI-driven tools into a practical healthcare management workflow.
\end{abstract}

\begin{IEEEkeywords}
Healthcare, Artificial Intelligence, Machine Learning, Django, React, Medical Triage, Disease Prediction
\end{IEEEkeywords}

\section{Introduction}
The healthcare industry faces continuous challenges regarding resource allocation, wait times, and the efficiency of patient triage. With the advent of scalable web technologies and accessible artificial intelligence (AI), there is a significant opportunity to automate the preliminary phases of patient care. 

MedPredictX was conceptualized to address these challenges by providing a digital-first approach to medical triage. It serves as an intelligent intermediary that evaluates patient symptoms, recommends appropriate medical specialists, and even predicts potential diseases using historical medical data. By providing patients with immediate insights into their health conditions, MedPredictX reduces unnecessary anxiety and helps prioritize severe cases.

The core objectives of MedPredictX are:
\begin{itemize}
    \item To provide an intuitive user interface for patients to log symptoms and book appointments.
    \item To integrate a high-speed AI recommender system using the Groq API.
    \item To deploy a local machine learning predictor that analyzes 132 symptom points for accurate disease forecasting.
    \item To maintain a secure, scalable backend for managing patient profiles, medical records, and doctors' schedules.
\end{itemize}

\section{Literature Review}
The integration of Machine Learning (ML) in healthcare has been extensively researched. Previous systems have utilized Decision Trees, Support Vector Machines (SVM), and Neural Networks to diagnose diseases ranging from diabetes to cardiovascular conditions.

While traditional telemedicine platforms focus primarily on video conferencing and basic appointment scheduling, modern systems increasingly rely on Natural Language Processing (NLP) to parse patient symptoms. However, many of these systems rely entirely on cloud-based LLMs, which can be subject to rate limits and latency. MedPredictX innovates by utilizing a hybrid approach: employing a high-speed LLM (Groq) for conversational triage and a deterministic, localized Random Forest model for pure symptom-to-disease classification.

\section{System Architecture}
MedPredictX follows a classic three-tier architecture comprising a presentation tier, an application tier, and a data tier. 

\subsection{High-Level Overview}
The system is divided into two primary repositories: the frontend (React/Vite) and the backend (Django REST Framework). Communication between the two occurs via RESTful APIs over HTTP, secured by JSON Web Tokens (JWT).

% PLACEHOLDER INSTRUCTION: Create a block diagram showing the Frontend (React), Backend (Django), Database (SQLite/PostgreSQL), and the two AI components (Groq API and Local ML Model).
\begin{figure}[htbp]
\centerline{\includegraphics[width=0.8\columnwidth]{placeholder_architecture.png}}
\caption{System Architecture Diagram. (PLACEHOLDER: Insert a flowchart showing React communicating with Django REST API, which in turn interacts with the local ML Joblib model and the external Groq API.)}
\label{fig:architecture}
\end{figure}

\subsection{Database Design}
The data tier manages relational entities such as Users, Doctors, Appointments, and Medical Records. The Django ORM handles schema migrations and ensures data integrity through relational constraints.

\section{Frontend Development and User Interface}
The presentation tier is built using React, providing a dynamic and responsive Single Page Application (SPA). The frontend is responsible for capturing user input, visualizing data, and handling local state management.

\subsection{User Interface Components}
The user interface is designed with a focus on accessibility and ease of use. Key components include:
\begin{itemize}
    \item \textbf{Dashboard}: A central hub for users to view upcoming appointments and recent medical records.
    \item \textbf{Symptom Predictor}: A specialized form allowing users to select from 132 specific symptoms.
    \item \textbf{AI Recommender}: A chat-like interface where users can describe their issues in natural language.
\end{itemize}

% PLACEHOLDER INSTRUCTION: Capture a screenshot of the main dashboard showing the navigation bar and user summary.
\begin{figure}[htbp]
\centerline{\includegraphics[width=0.8\columnwidth]{placeholder_dashboard.png}}
\caption{User Dashboard Interface. (PLACEHOLDER: Insert a screenshot of the React frontend showing the Patient Dashboard.)}
\label{fig:dashboard}
\end{figure}

\subsection{State Management and API Integration}
React hooks (\texttt{useState}, \texttt{useEffect}) manage the component state. API calls are handled via \texttt{axios}, which automatically attaches the JWT bearer token to outbound requests for authenticated endpoints.

\section{Backend Implementation and API Design}
The application tier is powered by Django and the Django REST Framework (DRF). This tier enforces business logic, handles AI model integration, and serves data to the frontend.

\subsection{Authentication and Authorization}
Security is paramount in healthcare applications. MedPredictX utilizes \texttt{rest\_framework\_simplejwt} to issue access and refresh tokens. All sensitive endpoints are protected by \texttt{IsAuthenticated} permission classes, ensuring that medical records and appointments are only accessible to their respective owners.

\subsection{RESTful Endpoints}
The backend exposes several critical endpoints:
\begin{itemize}
    \item \texttt{/api/auth/*}: Registration, login, and token refresh.
    \item \texttt{/api/appointments/}: CRUD operations for doctor appointments.
    \item \texttt{/api/medical-records/}: Management of patient history.
\end{itemize}

% PLACEHOLDER INSTRUCTION: Capture a screenshot of the Django REST Framework browsable API or a Postman request showing a successful API response.
\begin{figure}[htbp]
\centerline{\includegraphics[width=0.8\columnwidth]{placeholder_api.png}}
\caption{Backend API Response. (PLACEHOLDER: Insert a screenshot of an API tool like Postman demonstrating a successful JSON response from the backend.)}
\label{fig:api}
\end{figure}

\section{Artificial Intelligence and Machine Learning Integration}
The defining feature of MedPredictX is its dual-layered artificial intelligence engine.

\subsection{AI Doctor Recommender (Groq API)}
The recommender system utilizes the Groq API (running the LLaMA model) to parse natural language descriptions of patient ailments. It returns structured JSON data containing:
\begin{itemize}
    \item The recommended medical specialty.
    \item An urgency level (e.g., Low, Medium, High).
    \item Immediate actions the patient should take.
\end{itemize}
This component acts as an intelligent triage nurse, directing the patient to the correct specialist within the application.

\subsection{Local Disease Predictor (Random Forest)}
To provide concrete disease predictions without relying solely on external APIs, MedPredictX features a local Machine Learning model. 
A \texttt{RandomForestClassifier} was trained on a dataset (\texttt{Training.csv}) comprising 4,920 records mapped across 132 distinct binary symptom features (e.g., itching, skin rash, continuous sneezing) against a target prognosis (e.g., Fungal infection, Allergy).

The training script processes the CSV, extracts the feature names, trains the ensemble model, and serializes the state using \texttt{joblib}. When the frontend submits an array of symptoms to the \texttt{/api/predict/} endpoint, the backend constructs a 132-dimensional feature vector, invokes the model, and returns the predicted disease in milliseconds.

% PLACEHOLDER INSTRUCTION: Capture a screenshot of the Predictor interface where multiple symptoms are selected and a disease prediction is shown.
\begin{figure}[htbp]
\centerline{\includegraphics[width=0.8\columnwidth]{placeholder_predictor.png}}
\caption{ML Disease Predictor. (PLACEHOLDER: Insert a screenshot of the frontend Predictor component showing selected symptoms and the resulting ML prediction.)}
\label{fig:predictor}
\end{figure}

\section{Results and Discussion}
The implemented system successfully achieves its primary objectives. The separation of the AI Recommender and the Local Predictor provides a robust fail-safe. 

\subsection{Model Accuracy}
The Random Forest model, tested on a localized subset of the symptom dataset, demonstrates near-perfect classification accuracy due to the highly deterministic nature of the training data. The vectorization approach ensures that the backend can process requests with minimal computational overhead.

\subsection{System Performance}
API response times for standard CRUD operations average below 100ms on a local environment. The Groq API Integration provides responses typically within 1-2 seconds, while the local ML predictor executes in under 50ms, highlighting the efficiency of utilizing serialized \texttt{joblib} models for production inference.

% PLACEHOLDER INSTRUCTION: Create a graph or table showing the response times of the Groq API vs the Local ML Predictor.
\begin{figure}[htbp]
\centerline{\includegraphics[width=0.8\columnwidth]{placeholder_performance.png}}
\caption{Performance Comparison. (PLACEHOLDER: Insert a bar chart comparing the inference time of the external Groq API versus the local Random Forest model.)}
\label{fig:performance}
\end{figure}

\section{Conclusion and Future Work}
MedPredictX demonstrates the viability of integrating both local machine learning models and cloud-based large language models into a unified healthcare web application. The platform successfully manages patient data while providing highly accurate triage recommendations and disease predictions.

\subsection{Future Scope}
Future enhancements for MedPredictX include:
\begin{itemize}
    \item \textbf{Hyperparameter Tuning}: Further optimization of the local ML model to handle noisy or incomplete symptom data.
    \item \textbf{Telemedicine Integration}: Adding WebRTC capabilities for live video consultations with the recommended doctors.
    \item \textbf{Continuous Learning}: Implementing a feedback loop where doctor-verified diagnoses are fed back into the training dataset to continuously improve the local model's accuracy.
\end{itemize}

\section*{Acknowledgment}
The authors would like to thank the developers of the open-source libraries utilized in this project, including the Django Software Foundation, React, Scikit-Learn, and the Groq team for their API services.

\begin{thebibliography}{00}
\bibitem{b1} Django Software Foundation, ``Django Documentation,'' https://docs.djangoproject.com/
\bibitem{b2} React, ``A JavaScript library for building user interfaces,'' https://reactjs.org/
\bibitem{b3} F. Pedregosa et al., ``Scikit-learn: Machine Learning in Python,'' JMLR 12, pp. 2825-2830, 2011.
\bibitem{b4} Groq API Documentation, ``Fast AI Inference,'' https://console.groq.com/docs
\end{thebibliography}

\end{document}
